\part{Quiver / Spin-Chain Models}

\section{Bethe ansatz}

\begin{exercise}[The inhomogeneous rational spin chain and a one-root Bethe state]
Let $V=\C^{\{+,-\}}$ be the arrow space of the section
Row-to-row transfer matrices. For $\epsilon\ne0$ use
$R(u)=u\id+\epsilon P$ and
$\mathcal T_0(u)=R_{0N}(u-\theta_N)\cdots R_{01}(u-\theta_1)$.
With the auxiliary kernels $A,B,C,D$ of the section Commuting
families and nonzero $\kappa_\pm$, put
$t_\kappa(u)=\kappa_+A(u)+\kappa_-D(u)$.
Let $\Omega$ be the all-up configuration.
\begin{enumerate}[label=(\alph*)]
\item Prove $[t_\kappa(u),t_\kappa(v)]=0$ by RTT and the
commutation of $R$ with the product of the two diagonal twists.
Compute
$A(u)\Omega=a(u)\Omega$, $D(u)\Omega=d(u)\Omega$, $C(u)\Omega=0$,
where $a(u)=\prod_j(u-\theta_j+\epsilon)$ and
$d(u)=\prod_j(u-\theta_j)$.
\item Extract the two exchange relations involving $A,B$ and
$D,B$ from RTT. For a nonzero state $B(\lambda)\Omega$ derive
\[
 \Lambda(u)=\kappa_+a(u)\frac{u-\lambda-\epsilon}{u-\lambda}
          +\kappa_-d(u)\frac{u-\lambda+\epsilon}{u-\lambda},
 \qquad
 \kappa_+a(\lambda)=\kappa_-d(\lambda).
\]
Show that the Bethe equation cancels the apparent pole at
$u=\lambda$ and the unwanted vector $B(u)\Omega$.
\end{enumerate}
\end{exercise}

\section{Quiver varieties}

\begin{exercise}[A framed one-vertex quiver and the cotangent Grassmannian]
A quiver is a finite directed graph. Consider the quiver with one
vertex, no arrows, a framing space $W$, and a vertex space $U$,
both finite-dimensional. Its doubled framed data are
$I:W\to U$, $J:U\to W$, with action
$g\cdot(I,J)=(gI,Jg^{-1})$ of $\GL(U)$.
Use the trace pairing to equip this space with the symplectic
form $\omega((I_1,J_1),(I_2,J_2))=\Tr(I_1J_2-I_2J_1)$.
The moment map is $\mu(I,J)=IJ\in\End(U)$; stability means that
$I$ is surjective. The quiver variety is
\[
 \mathfrak M(r,W)=\{(I,J):IJ=0,\ I\text{ surjective},
                       \ \dim U=r\}/\GL(U).
\]
\begin{enumerate}[label=(\alph*)]
\item For $\xi\in\End(U)$ verify
$\iota_{\xi^\#}\omega=\dd\Tr(\xi\mu)$, fixing this sign
convention for moment maps.
\item Put $S=\ker I$. Prove that $J$ has image in $S$ and
descends to a map $W/S\to S$. Identify the orbit space with
the cotangent bundle of the Grassmannian of codimension-$r$
subspaces of $W$, using the trace duality between
$\Hom(S,W/S)$ and $\Hom(W/S,S)$.
\item For $r=1$, $\dim W=2$, identify $\mathfrak M$ with
$T^*\mathbb P(W)$ by recording the kernel line. Prove that
scaling the cotangent fibre retracts this space onto
$\mathbb P(W)$. Compute the varieties for $r=0$ and $r=2$.
\end{enumerate}
\end{exercise}

\section{\texorpdfstring{Cohomology and $K$-theory}{Cohomology and K-theory}}

\begin{exercise}[The cohomology of a projective line and its line bundles]
Use the cellular cohomology and bundle group completion of the
sections Higher gauge theory and Generalized cohomology.
The first Chern class of a complex line bundle $L$ is the pullback
of the generator of $H^2(\mathbb P(\C^\infty);\Z)$ under its
classifying map, normalized so that the dual tautological line
on $\mathbb P(\C^2)$ has integral one.
\begin{enumerate}[label=(\alph*)]
\item Build $\mathbb P(\C^n)$ by its cells of dimensions
$0,2,\dots,2n-2$. Intersect transverse projective hyperplanes
to show that the powers of their degree-two class give
$H^*(\mathbb P(\C^n);\Z)=\Z[x]/(x^n)$.
Deduce $H^*(T^*\mathbb P(W);\Z)=\Z[x]/(x^2)$ for $\dim W=2$.
\item Let $L$ be the tautological line on $\mathbb P(W)$.
Use $0\to L\to W\to W/L\to0$ and its determinant to prove
$[L]+[L]^{-1}=2$ when $W$ is two-dimensional and constant.
Deduce $([L]-1)^2=0$ in $K^0(\mathbb P(W))$ and calculate
the rank and first Chern class of $[L]-1$.
\end{enumerate}
\end{exercise}

\begin{exercise}[Equivariant cohomology of the two-site quiver variety]
Specify the physical framing decomposition $W=W_1\oplus W_2$
into lines and let a torus scale the two lines independently.
Let an additional circle scale cotangent fibres.
Equivariant cohomology is
$H_T^*(X;\C)=H^*(ET\times_T X;\C)$, where $ET$ is a
contractible free $T$-space; for circles use the union of odd
unit spheres in $\C^\infty$.
Write $a_1,a_2,\epsilon$ for the degree-two character classes,
$B=\C[a_1,a_2,\epsilon]$, and $t=c_1(L)$.
\begin{enumerate}[label=(\alph*)]
\item Use finite projective-space approximations to $BT$ and
the two-cell calculation of the preceding exercise to prove
\[
 H_T^*(T^*\mathbb P(W);\C)
 =B[t]/((t-a_1)(t-a_2)).
\]
Derive the relation from the tautological line and quotient
line whose sum is the pulled-back framing bundle.
\item Restrict to the two fixed points $p_i=\mathbb P(W_i)$
and prove $t|_{p_i}=a_i$. Show that restriction is injective
and that a pair $(f_1,f_2)\in B^2$ is in its image exactly
when $f_1-f_2$ is divisible by $a_1-a_2$.
\end{enumerate}
\end{exercise}

\section{Stable envelopes}

\begin{exercise}[Two chambers and the stable-envelope maps of a projective line]
For the two fixed points in the preceding exercise choose opposite
orders $p_1<p_2$ and $p_2<p_1$, called chambers $+$ and $-$.
Fix the polarization and signs by the following normalization:
the stable-envelope class for $p_1$ in chamber $+$ vanishes
at $p_2$ and equals $a_1-a_2$ at $p_1$; the class for $p_2$
has restrictions $\epsilon$ and $a_2-a_1+\epsilon$.
In chamber $-$ the classes have restrictions
$(a_1-a_2+\epsilon,\epsilon)$ and $(0,a_2-a_1)$.
These support, normalization, and degree conditions define the
stable envelopes in this two-fixed-point model.
\begin{enumerate}[label=(\alph*)]
\item Use the divisibility criterion to prove existence and
uniqueness, and obtain
\[
 \operatorname{Stab}_+(p_1)=t-a_2,\quad
 \operatorname{Stab}_+(p_2)=t-a_1+\epsilon,\quad
 \operatorname{Stab}_-(p_1)=t-a_2+\epsilon,\quad
 \operatorname{Stab}_-(p_2)=t-a_1 .
\]
\item Put $u=a_1-a_2$ and invert $u$, $u-\epsilon$,
$u+\epsilon$. Compute the chamber-change operator
$\operatorname{Stab}_-^{-1}\operatorname{Stab}_+$ on the
fixed-point function space:
\[
 \frac1{u+\epsilon}
 \begin{pmatrix}u&\epsilon\\ \epsilon&u\end{pmatrix}.
\]
The entries here are indexed by the two geometrically fixed points.
\item Add the two one-point quiver varieties $r=0,2$, with
chamber change equal to one. Identify the resulting operator
with $(u\id+\epsilon P)/(u+\epsilon)$ on two arrow spaces.
Verify its Yang--Baxter identity using the rational limit of
the section Spectral parameter.
Compare this concrete construction with \cite{maulik-okounkov}.
\end{enumerate}
\end{exercise}

\section{Geometric representation theory}

\begin{exercise}[Permutation symmetry and the two-site representation]
Let $\mathcal W=\bigoplus_{r=0}^2
H_T^*(\mathfrak M(r,W))$ with the character differences inverted,
and identify its fixed-point function space with the four spin
configurations of $V^{\otimes2}$ by the occupied subset of
$\{W_1,W_2\}$. On each arrow space use the spin operators of
the section Yang--Baxter equation.
\begin{enumerate}[label=(\alph*)]
\item Show that the raising and lowering operators on the
occupied subsets, defined by removing or adding one label
with coefficient one, satisfy
$[S^z_{\rm tot},S^\pm_{\rm tot}]=\pm S^\pm_{\rm tot}$ and
$[S^+_{\rm tot},S^-_{\rm tot}]=2S^z_{\rm tot}$.
\item Prove that chamber change commutes with this action.
Use the intrinsic projections $(\id\pm P)/2$ to split the
two-site space into symmetric and alternating subspaces.
Compute their dimensions and the chamber-change eigenvalues
$1$ and $(u-\epsilon)/(u+\epsilon)$.
\item For $\theta_1=\theta_2=0$, $\kappa_+=\kappa_-=1$,
solve the one-root Bethe equation and show that
$B(-\epsilon/2)\Omega$ belongs to the alternating line.
Identify the complementary spin-one multiplet by the total
spin action.
\end{enumerate}
\end{exercise}

\begin{exercise}[Prediction: the exchange splitting of a two-spin bond]
Let two physical spins of magnitude $\tfrac12$ interact through
$H=J(\mathbf S_1\cdot\mathbf S_2-\tfrac14)$ with $J>0$
and $\hbar=1$ in the spin commutation relations.
Use the quiver chamber-change operator
$\mathcal R(u)=(u\id+\epsilon P)/(u+\epsilon)$.
\begin{enumerate}[label=(\alph*)]
\item Prove
$H=\frac{J\epsilon}{2}\mathcal R(0)^{-1}\mathcal R'(0)
=\frac J2(P-\id)$.
Derive a singlet energy $-J$ and a triplet energy $0$.
\item In a magnetic field with Zeeman term
$-hS^z_{\rm tot}$ calculate all four energies and the first
ground-state crossing at $h=J$.
\item Derive the zero-field partition function
$Z=e^{\beta J}+3$ and susceptibility to $h$,
$\chi=2\beta/(e^{\beta J}+3)$.
State the singlet-triplet spectroscopic gap $J$ and specify
that isotropic exchange, two spin-$\tfrac12$ degrees of freedom,
and negligible coupling to other bonds are the model assumptions.
\end{enumerate}
\end{exercise}
